\input{../tex-inputs/latex-leseansicht-vorspann}

\section[Arthur Schnitzler an Gustav Schwarzkopf, 25. 1. 1914]{L04503 Arthur Schnitzler an Gustav Schwarzkopf, 25. 1. 1914}
\nopagebreak\mylabel{L04503v}
\rehead{ }\normalsize\beginnumbering\briefempfaengerindex{Schwarzkopf, Gustav@\textsc{Schwarzkopf, Gustav}!zzzSchnitzler, Arthur@\emph{von Arthur Schnitzler}!1914-01-251@{25. 1. 1914}|(be}
\toendnotes[C]{\smallbreak\pagebreak[2]}
\correspDesc{Versand  durch Arthur Schnitzler am 25. 1. 1914 in Regensburg
\newline{}Übermittlung  am 25. 1. 1914 in Regensburg
\newline{}Erhalt  durch Gustav Schwarzkopf im Zeitraum [26. 1. 1914 – 30. 1. 1914?] in Wien}\toendnotes[C]{\smallbreak}
\Standort{DLA, A:Schnitzler, HS.1985.1.1897.}
\physDesc{Bildpostkarte, 444 Zeichen
\newline{}Handschrift: Bleistift, deutsche Kurrent
\newline{}Versand: Stempel: »\nobreak{}\oindex{Regensburg@\textbf{Regensburg}|pwk}\textcolor{gray}{Regens}burg, 25. 1. 14, 3–4 N\nobreak{}«.  }\toendnotes[C]{\smallbreak}
\pstart
           \noindent{}\centering{}{\pb}\textcolor{gray}{\textbf{REGENSBURG Jacobstor}}\oindex{Jakobstor@\textbf{Jakobstor}|pw}\pend
           \pstart{}{\pb}Herrn \textsc{Gustav
                     Schwarzkopf}\pend{}\pstart{}Wien I\oindex{I., Innere Stadt@\textbf{I., Innere Stadt}|pw}\pend{}\pstart{}\textsc{Tiefer
                        Graben 17}\oindex{Wien@\textbf{Wien}!I., Innere Stadt@\textbf{I., Innere Stadt}!Tiefer Graben 17@\textbf{Tiefer Graben 17}|pw}\pend{}{\bigskip}\vspace{1em}
\pstart
           \noindent{}Lieber Guſtav, gern hätt ich Ihnen  noch Adieu geſagt; aber leider
               wars mir nicht mehr möglich{ }ſtadtwärts zu{ }ſpaziren. Wir{ }ſehn uns hoffentlich bald
               wieder – da wir bei der Mordskälte wohl auf Berlin\oindex{Berlin@\textbf{Berlin}|pw} verzichten werden\pend
           
\pstart
           Bitte um ein Wort nach München\oindex{München@\textbf{München}|pw} (Hotel Marienbad\oindex{Hotel Marienbad [München]@\textbf{Hotel Marienbad [München]}|pw}) {\pb}wies
               Ihnen geht. – \label{K_L04503-1v}\edtext{Das geſtrige Concert\eventindex{Hotel Karmeliten@\textbf{Hotel Karmeliten}!Konzert Olga Schnitzler und Vera Schapira, 24.1.1914@Konzert Olga Schnitzler und Vera Schapira, 24.1.1914|pw}}{\lemma{\textnormal{\emph{Das gestrige Concert}}}\Cendnote{\textnormal{Konzert Olga Schnitzler und Vera Schapira, 24.1.1914. In: A. S.: \emph{Kulturveranstaltungen}, \url{https://schnitzler-kultur.acdh.oeaw.ac.at/pmb183645.html}.}}}\label{K_L04503-1}
               iſt recht gut verlaufen – allerlei Heiteres wollen wir Ihnen mündlich
                  erzähl\textcolor{gray}{en}\pend
           
\pstart
           Herzlichſt{\\[\baselineskip]}Ihr \spacefill\mbox{A.}\pend
           \leftskip=0em{}
\pstart
           \raggedleft{}25/1 914.\pend
           \selectlanguage{ngerman}\endnumbering\briefempfaengerindex{Schwarzkopf, Gustav@\textsc{Schwarzkopf, Gustav}!zzzSchnitzler, Arthur@\emph{von Arthur Schnitzler}!1914-01-251@{25. 1. 1914}|)be}\mylabel{L04503h}
\begin{anhang}
\end{anhang}\newcommand{\dateiname}{L04503}\newcommand{\titel}{Arthur Schnitzler an Gustav Schwarzkopf, 25. 1. 1914}\newcommand{\editorInnen}{Herausgegeben von Jahnke, SelmaMüller, Martin Anton}\input{../tex-inputs/latex-leseansicht-abspann}
